\chapter{Implementación}
\label{chap:implementacion}

\ph{Este capítulo describe \textbf{cómo} se ha materializado el diseño en código:
qué tecnologías se han usado, qué módulos se han implementado y qué decisiones
concretas se han tomado en cada uno.}

\section{Análisis de tecnologías}
\label{sec:analisis_tecnologias}

\ph{Justifica la elección de las tecnologías principales del proyecto.
Si has comparado varias opciones, hazlo explícito (es más convincente que
afirmar la elegida sin alternativas).}

\subsection*{\ph{Necesidad técnica 1}}

\ph{Plantea el problema concreto que motivó la elección
(p.\,ej., procesar archivos en formato X, validar reglas, automatizar Y\dots)
y qué características debe tener la tecnología elegida.}

\subsection*{\ph{Necesidad técnica 2}}

\ph{Otra necesidad que orienta una decisión tecnológica.}

\subsection*{\ph{Necesidad técnica 3}}

\ph{Otra necesidad que orienta una decisión tecnológica
(p.\,ej., experiencia de usuario, despliegue, rendimiento, mantenibilidad).}

La implementación del sistema requiere la utilización de tecnologías compatibles con
la arquitectura existente de UVLHub y con el ecosistema Flamapy. Por ello, la selección
de tecnologías se ha realizado teniendo en cuenta principalmente la integración con
la plataforma existente, la facilidad de desarrollo y la necesidad de disponer de un
entorno reproducible para la ejecución y despliegue del sistema.

\subsection*{Necesidad técnica 1: Integración con la plataforma existente}

El generador debe integrarse dentro de UVLHub, una plataforma web
existente desarrollada utilizando Flask. Por este motivo, la tecnología utilizada
debe permitir mantener la arquitectura actual basada en rutas, servicios y plantillas
Jinja2, evitando introducir cambios estructurales innecesarios en la aplicación.

\subsection*{Necesidad técnica 2: Desarrollo de la herramienta generadora}

La lógica de generación debe implementarse siguiendo la arquitectura de extensiones
proporcionada por Flamapy. Para ello, la tecnología utilizada debe permitir desarrollar
una herramienta independiente que encapsule la configuración del generador y la operación
encargada de producir modelos de características expresados en UVL.

\subsection*{Necesidad técnica 3: Ejecución y despliegue reproducible}

El sistema debe disponer de un entorno de ejecución reproducible tanto para el
desarrollo local como para facilitar las revisiones del proyecto. Por ello, es
necesario utilizar tecnologías que permitan contenerizar la aplicación, gestionar
sus dependencias y proporcionar una versión accesible del sistema desplegada en
un entorno remoto.

\subsubsection*{Comparativa de tecnologías del backend}

\ph{
\begin{itemize}
  \item \textbf{\ph{Tecnología A}}: \ph{descripción y ventajas/inconvenientes.}
  \item \textbf{\ph{Tecnología B}}: \ph{descripción y ventajas/inconvenientes.}
  \item \textbf{\ph{Tecnología C}}: \ph{descripción y ventajas/inconvenientes.}
\end{itemize}
}

\begin{itemize}

  \item \textbf{Flask:}
  Framework utilizado por UVLHub sobre el que se ha integrado el generador. Su principal
  ventaja es la compatibilidad directa con la arquitectura existente, permitiendo
  reutilizar las rutas, servicios y mecanismos de configuración ya implementados.

  \item \textbf{Django:}
  Framework web completo con numerosas funcionalidades integradas, como ORM o gestión
  de usuarios. Sin embargo, su utilización habría requerido adaptar la arquitectura
  existente de UVLHub, introduciendo una complejidad innecesaria.

  \item \textbf{FastAPI:}
  Framework orientado principalmente al desarrollo de APIs de alto rendimiento. Aunque
  ofrece ventajas para servicios desacoplados, no resulta adecuado en este caso al no
  encajar con la arquitectura basada en Flask utilizada por UVLHub.

\end{itemize}

\subsubsection*{Comparativa de tecnologías del frontend}

\ph{
\begin{itemize}
  \item \textbf{\ph{Tecnología A}}: \ph{descripción y ventajas/inconvenientes.}
  \item \textbf{\ph{Tecnología B}}: \ph{descripción y ventajas/inconvenientes.}
  \item \textbf{\ph{Tecnología C}}: \ph{descripción y ventajas/inconvenientes.}
\end{itemize}
}

\begin{itemize}

\item \textbf{Jinja2, JavaScript y Bootstrap:}
Tecnologías utilizadas por UVLHub para la construcción de la interfaz web. En
particular, Jinja2 constituye el motor de plantillas integrado con Flask,
permitiendo generar las vistas de la aplicación utilizando una arquitectura
compatible con el backend existente. Esta opción permite incorporar el flujo
tipo \textit{wizard} desarrollado para el generador manteniendo la coherencia
visual y estructural con el resto de la plataforma.

\item \textbf{React:}
Biblioteca ampliamente utilizada para el desarrollo de interfaces dinámicas
basadas en componentes reutilizables. Sin embargo, su utilización habría
requerido introducir una arquitectura frontend independiente de la empleada
actualmente por UVLHub, aumentando la complejidad de integración del generador.

\item \textbf{Vue.js:}
Framework orientado al desarrollo de interfaces reactivas mediante componentes.
Aunque ofrece una solución ligera para construir interfaces modernas, no fue
seleccionado debido a que no forma parte de la arquitectura frontend existente
en UVLHub.

\end{itemize}

\subsubsection*{Análisis de la base de datos}

\ph{
\begin{enumerate}
  \item \textbf{\ph{Opción A}}: \ph{justificación.}
  \item \textbf{\ph{Opción B}}: \ph{justificación.}
\end{enumerate}
}

El sistema desarrollado no requiere una base de datos propia, ya que la
naturaleza del generador no implica el almacenamiento permanente de
configuraciones ni de modelos generados. La información necesaria durante el
flujo de configuración se mantiene temporalmente mediante la \textit{session}
de Flask, permitiendo ejecutar el proceso completo sin introducir una capa de
persistencia adicional, lo que hace innecesario realizar aquí una comparación.

\subsubsection*{Plataformas de despliegue}

\ph{
\begin{itemize}
  \item \textbf{\ph{Plataforma A}}: \ph{descripción y casuística.}
  \item \textbf{\ph{Plataforma B}}: \ph{descripción y casuística.}
  \item \textbf{\ph{Plataforma C}}: \ph{descripción y casuística.}
\end{itemize}
}

\begin{itemize}

  \item \textbf{Render:}
  Plataforma utilizada para desplegar una versión accesible de UVLHub durante las
  revisiones del proyecto. Su elección estuvo motivada tanto por la facilidad de
  configuración e integración con el repositorio del proyecto como por la
  experiencia previa personal con esta plataforma, evitando la
  necesidad de invertir tiempo adicional en el aprendizaje de una nueva
  tecnología de despliegue.

  \item \textbf{Servicios cloud tradicionales (AWS, Azure, Google Cloud):}
  Ofrecen una mayor capacidad de configuración y escalabilidad, aunque requieren
  una gestión más compleja de infraestructura. Debido a que el objetivo del
  despliegue era facilitar la revisión funcional del sistema y no proporcionar
  un servicio preparado para producción, no resultaban necesarios para este
  proyecto.

  \item \textbf{Heroku:}
  Plataforma orientada al despliegue sencillo de aplicaciones web con una
  filosofía similar a Render. Aunque habría permitido obtener una solución
  equivalente, no fue seleccionada al disponer ya de experiencia previa y una
  configuración establecida con Render.

\end{itemize}

\subsubsection*{Tecnologías finales seleccionadas}

\ph{
\begin{itemize}
  \item \textbf{Desarrollo}: \ph{tecnologías elegidas para backend / frontend.}
  \item \textbf{Base de datos}: \ph{motor elegido.}
  \item \textbf{Contenedor}: \ph{tecnología de contenerización (Docker, Podman\dots).}
  \item \textbf{Despliegue}: \ph{plataforma final.}
\end{itemize}
}

\begin{itemize}
    \item \textbf{Desarrollo backend:} Python, Flask y Flamapy, utilizando la
    estructura base para la creación de extensiones proporcionada por este último.

    \item \textbf{Desarrollo frontend:} Plantillas Jinja2, JavaScript y Bootstrap,
    manteniendo la arquitectura visual existente de UVLHub.

    \item \textbf{Persistencia:} No aplica, al no requerirse almacenamiento
    permanente de configuraciones ni modelos generados.

    \item \textbf{Contenerización:} Docker, utilizado para construir y ejecutar
    el entorno completo de desarrollo local mediante contenedores.

    \item \textbf{Despliegue:} Render, utilizado para proporcionar una versión
    accesible del sistema durante las revisiones del proyecto.

\end{itemize}

\section{Tecnologías utilizadas}
\label{sec:tecnologias}

\ph{Tras la justificación, describe brevemente cómo se usa cada tecnología
en el proyecto. Una tecnología por párrafo, indicando módulo o archivo
en que aparece.}

\ph{Sugerencia de bloques:}

\ph{
\begin{itemize}
  \item \ph{Backend (framework, gestión de dependencias, configuración).}
  \item \ph{Autenticación (estrategia, almacenamiento de credenciales,
        gestión de sesión o token).}
  \item \ph{Persistencia (ORM, motor, migraciones).}
  \item \ph{Lógica de dominio (módulos en los que se divide).}
  \item \ph{Generación de artefactos / motor de plantillas (si aplica).}
  \item \ph{Capa de presentación (HTML/CSS/JS o framework).}
  \item \ph{Infraestructura de desarrollo (Docker, Compose, Make\dots).}
  \item \ph{Suite de pruebas (framework, marcadores, comandos).}
  \item \ph{Automatización CI/CD (proveedor y workflows principales).}
\end{itemize}
}

Tras la selección de tecnologías realizada en el apartado anterior, esta sección
describe cómo se han utilizado durante la implementación del sistema, indicando
los principales componentes en los que aparecen.

\subsection*{Backend y lógica del servidor}

La implementación del backend de UVLHub está desarrollada utilizando Python y
Flask. Dentro del módulo \textit{generator} se encuentran las rutas encargadas
de gestionar el flujo de configuración del generador, implementadas en el
archivo \texttt{app/features/generator/routes.py}, así como los servicios
internos definidos en \texttt{app/features/generator/wizard.py}, responsables
de reconstruir y validar la configuración introducida por el usuario.

\begin{figure}[H]
    \centering
    \includegraphics[width=0.30\textwidth]{estructura_modulo_generator.png}
    \caption{Estructura del módulo generator de UVLHub. Elaboración propia}
    \label{fig:estructura_generator}
\end{figure}

Además, Flask proporciona el mecanismo de \textit{session}, utilizado durante el flujo de configuración para mantener temporalmente los parámetros del generador en forma de diccionario entre los diferentes pasos del \textit{wizard}. El servidor expone esta información en formato JSON mediante el endpoint \texttt{/generator/random/params-json}. Posteriormente, el navegador utiliza estos parámetros para construir la instancia de \texttt{FmgeneratorModel} mediante \texttt{from\_flat\_dict}, dentro del archivo \texttt{fmgen\_wrapper.py} ejecutado con Pyodide.


\subsection*{Generador compatible con Flamapy}

La lógica principal de generación se ha implementado como un componente compatible
con la arquitectura de extensiones proporcionada por Flamapy. Este componente
se encuentra dentro del directorio \texttt{fm\_generator} del proyecto UVLHub y
contiene el modelo de configuración del generador mediante la clase
\textit{FmgeneratorModel}, junto con las diferentes clases encargadas de
representar los parámetros específicos de generación.

La ejecución del proceso se encapsula en la operación
\texttt{GenerateFeatureModel}, ubicada dentro del módulo de operaciones
de la herramienta. Su método \texttt{execute} utiliza la configuración
definida en \texttt{FmgeneratorModel} para construir un objeto
\texttt{FeatureModel}, que se recupera mediante \texttt{get\_result}.
La conversión de este objeto a formato UVL se realiza posteriormente
mediante \texttt{UVLWriter}.

\begin{figure}[H]
    \centering
    \includegraphics[width=0.30\textwidth]{estructura_fm_generator.png}
    \caption{Estructura de la herramienta generadora basada en Flamapy. Elaboración propia}
    \label{fig:estructura_fm_generator}
\end{figure}


\subsection*{Capa de presentación}

La interfaz web del generador mantiene la arquitectura frontend existente en
UVLHub, utilizando plantillas Jinja2 integradas con Flask para la generación de
las vistas, junto con JavaScript y Bootstrap para gestionar la interacción y los
componentes visuales de la aplicación. Las diferentes pantallas que forman el
flujo tipo \textit{wizard} se encuentran definidas como plantillas HTML dentro
del directorio \texttt{app/features/generator/templates/generator}.

\begin{figure}[H]
    \centering
    \includegraphics[width=0.30\textwidth]{templates_generator.png}
    \caption{Plantillas HTML utilizadas para los pasos del \textit{wizard}. Elaboración propia}
    \label{fig:templates_generator}
\end{figure}


\subsection*{Ejecución mediante Pyodide y WebAssembly}

Para permitir la ejecución del generador directamente en el navegador del
usuario se ha utilizado Pyodide, un entorno basado en WebAssembly que permite
ejecutar código Python en el cliente. Los scripts encargados de configurar
Pyodide y cargar las dependencias necesarias para la generación se encuentran
en el directorio \texttt{app/features/generator/assets/js}, donde también se
incluyen los paquetes \textit{wheel} necesarios para la ejecución de la herramienta
generadora.

\begin{figure}[H]
    \centering
    \includegraphics[width=0.30\textwidth]{pyodide_assets.png}
    \caption{Scripts y dependencias utilizadas para la ejecución con Pyodide. Elaboración propia}
    \label{fig:pyodide_assets}
\end{figure}


\subsection*{Infraestructura de desarrollo}

Durante el desarrollo local se ha utilizado Docker junto con Docker Compose para
definir y ejecutar el entorno completo necesario para UVLHub. Los ficheros de
configuración asociados se encuentran en el directorio
\texttt{docker}, donde se incluyen las definiciones necesarias tanto para los
entornos de desarrollo como para los entornos de producción.

\begin{figure}[H]
    \centering
    \includegraphics[width=0.30\textwidth]{docker_structure.png}
    \caption{Configuración de Docker en UVLHub. Elaboración propia}
    \label{fig:docker_structure}
\end{figure}

Aunque el entorno contiene diferentes servicios auxiliares necesarios para el
funcionamiento global de UVLHub, el único servicio con una relación directa con
el generador es Redis. Este servicio se utiliza durante el despliegue de
revisiones para almacenar temporalmente la información asociada a la
\textit{session} de Flask, permitiendo mantener la configuración actual del
usuario durante el flujo de generación.


\subsection*{Suite de pruebas}

Las pruebas del sistema se
ejecutan mediante \textit{Rosemary}, una CLI (\textit{Command Line Interface}) incluida dentro del entorno de UVLHub \cite{rosemary}.

La suite de pruebas incluye comprobaciones sobre diferentes componentes del
sistema, incluyendo pruebas de flujo del \textit{wizard}, validación de
parámetros y pruebas de interfaz mediante Selenium, permitiendo verificar el
correcto funcionamiento de la interacción del usuario con la aplicación.

\begin{figure}[H]
    \centering
    \includegraphics[width=0.30\textwidth]{tests_generator.png}
    \caption{Suite de pruebas del módulo generator. Elaboración propia}
    \label{fig:tests_generator}
\end{figure}


\subsection*{Automatización CI/CD}

Para automatizar las tareas de integración y despliegue se ha utilizado GitHub
Actions. Los flujos definidos permiten realizar comprobaciones técnicas previas
sobre el código desarrollado antes de completar el despliegue de la aplicación.

Además, se ha configurado la automatización del despliegue, permitiendo
actualizar la versión desplegada del sistema de forma automática tras superar
los controles establecidos.

\section{Módulo de validación / lógica principal}
\label{sec:modulo_principal}

\ph{Describe el módulo central de la lógica de negocio. Suele ser el que aplica
las reglas del dominio (validador, motor de reglas, planificador\dots).}

\ph{Cubre:}
\ph{
\begin{itemize}
  \item \ph{Función o clase principal y su contrato (entrada/salida).}
  \item \ph{Estructura del resultado (campos, tipos, semántica).}
  \item \ph{Cómo lo consumen otros módulos (API, generadores\dots).}
\end{itemize}
}

El módulo de validación se estructura en dos niveles diferenciados. Por un lado, el servicio 
\texttt{GeneratorWizardService} del módulo \texttt{generator} realiza una primera validación durante el proceso de configuración, comprobando que los valores introducidos por el usuario cumplen las restricciones básicas necesarias para continuar el flujo.

Por otro lado, el modelo de dominio \texttt{FmgeneratorModel} incorpora la validación semántica de la configuración completa del generador, comprobando la coherencia entre sus diferentes parámetros antes de iniciar la generación.

La generación de modelos de características se encapsula en la
operación \texttt{GenerateFeatureModel}, que recibe una instancia de
\texttt{FmgeneratorModel} y genera un único objeto \texttt{FeatureModel}
por ejecución. El método \texttt{execute} devuelve la propia operación,
mientras que \texttt{get\_result} permite recuperar el modelo generado.
Antes de ejecutar la generación, el método \texttt{execute} realiza
una última validación mediante:

\begin{lstlisting}[language=Python]
self.model.validate()
\end{lstlisting}

garantizando que únicamente configuraciones consistentes alcanzan la fase de generación.

\subsection{Validación sintáctica / comprobaciones básicas}
\label{subsec:validacion_sintactica}

La validación sintáctica se encuentra implementada principalmente en el fichero 
\texttt{wizard.py}, mediante un conjunto de funciones asociadas a cada paso del asistente de configuración:

\begin{lstlisting}[language=Python]
validate_step1_form()
validate_step2_form()
validate_step3_form()
validate_step4_form()
validate_step5_form()
validate_step6_form()
\end{lstlisting}

Estas funciones comprueban que los valores introducidos por el usuario cumplen las restricciones básicas necesarias antes de realizar la validación semántica del modelo de dominio.

\subsubsection*{Validación de parámetros generales}

Durante el primer paso del asistente se validan los parámetros básicos de generación. El número de modelos y la semilla deben cumplir:

\[
1 \leq NUM\_MODELS \leq 1000
\]

\[
SEED \in \mathbb{Z}^{+}
\]



\subsubsection*{Validación de rangos numéricos}

Los parámetros asociados a cantidades, límites y tamaños deben mantenerse dentro de rangos válidos. De forma general:

\[
VALUE \geq 1
\]

Además, todos los pares de límites mínimo-máximo cumplen:

\[
MIN\_VALUE \leq MAX\_VALUE
\]

Esta restricción se aplica a parámetros como número de features, número de restricciones, número de atributos o rangos definidos para atributos manuales.



\subsubsection*{Validación de distribuciones probabilísticas}

Las configuraciones probabilísticas utilizadas durante la generación deben representar distribuciones válidas. Para cualquier conjunto de probabilidades asociado a una decisión aleatoria:

\[
\forall p_i \in P,\quad 0 \leq p_i \leq 1
\]

\[
\sum_{p_i \in P} p_i = 1
\]

Estas comprobaciones se aplican a las distribuciones de relaciones entre features, operadores booleanos, operaciones aritméticas, tipos de restricciones y tipos de atributos.



\subsubsection*{Validación de dependencias básicas}

Las opciones del asistente presentan dependencias entre niveles de generación que deben cumplirse antes de continuar con la configuración:

\[
FEATURE\_CARDINALITY \Rightarrow ARITHMETIC\_LEVEL
\]

\[
AGGREGATE\_FUNCTIONS \Rightarrow ARITHMETIC\_LEVEL
\]

\[
STRING\_CONSTRAINTS \Rightarrow TYPE\_LEVEL
\]



\subsubsection*{Validación de atributos manuales}

En el modo de definición manual, cada atributo debe cumplir las restricciones básicas de configuración:

\[
name(Attribute) \neq \emptyset
\]

\[
0 \leq attach\_probability(Attribute) \leq 1
\]

Para atributos con dominio definido mediante rangos:

\[
0 \leq MIN\_ATTRIBUTE\_VALUE \leq MAX\_ATTRIBUTE\_VALUE
\]

Además, la posibilidad de utilizar un atributo dentro de una restricción depende de los niveles activos del generador:

\[
UseInConstraints(Attribute_{Integer|Real})
\Rightarrow Arithmetic\_Level
\]

\[
UseInConstraints(Attribute_{String})
\Rightarrow Type\_Level \land String\_Constraints
\]

De esta forma, la validación evita configuraciones incompatibles antes de que sean procesadas por la lógica de generación.

\subsection{Validación semántica y restricciones de dominio}
\label{subsec:validacion_semantica}

La validación semántica se encuentra encapsulada en el modelo 
\texttt{FmgeneratorModel}, definido dentro de 
\texttt{fm\_generator}. Este modelo agrupa la configuración completa del generador mediante diferentes bloques especializados:

\begin{table}[H]
\centering
\begin{tabular}{|c|c|}
\hline
\textbf{Configuración} & \textbf{Responsabilidad} \\
\hline
\texttt{NamingConfig} & Nombres y opciones de salida \\
\hline
\texttt{LevelConfig} & Niveles de expresividad UVL \\
\hline
\texttt{FeaturesConfig} & Parámetros de features \\
\hline
\texttt{HierarchyConfig} & Estructura jerárquica \\
\hline
\texttt{ConstraintsConfig} & Restricciones generadas \\
\hline
\texttt{AttributesConfig} & Gestión de atributos \\
\hline
\end{tabular}
\caption{Bloques de configuración de \texttt{FmgeneratorModel}}
\label{tab:generator_configs}
\end{table}

La validación completa se ejecuta mediante el método:

\begin{lstlisting}[language=Python]
FmgeneratorModel.validate()
\end{lstlisting}

que comprueba la coherencia global de la configuración antes de iniciar la generación.

En las comprobaciones de normalización, cada probabilidad debe
pertenecer al intervalo $[0,1]$. La implementación admite una
diferencia absoluta de hasta $10^{-6}$ entre la suma de las
probabilidades y uno, para contemplar las imprecisiones de la
representación numérica.

\subsubsection*{Coherencia de niveles UVL}

Los niveles de expresividad presentan dependencias jerárquicas entre ellos. Las reglas aplicadas son:

\[
BOOLEAN\_LEVEL = True
\]

\[
TYPE\_LEVEL \Rightarrow ARITHMETIC\_LEVEL
\]

\[
FEATURE\_CARDINALITY \Rightarrow ARITHMETIC\_LEVEL
\]

\[
AGGREGATE\_FUNCTIONS \Rightarrow ARITHMETIC\_LEVEL
\]

\[
STRING\_CONSTRAINTS \Rightarrow TYPE\_LEVEL
\]


\subsubsection*{Validación de la jerarquía de features}

La estructura jerárquica generada debe ser compatible con el número máximo de features:

\[
1 \leq MAX\_TREE\_DEPTH \leq MAX\_FEATURES
\]

Además, la selección probabilística de relaciones entre features requiere una distribución normalizada:

\[
P_{optional}+P_{mandatory}+P_{alternative}+P_{or}+P_{group}=1
\]


\subsubsection*{Validación de restricciones}

Las restricciones generadas dependen de los niveles activos. La distribución de tipos de restricciones se valida mediante:

\[
\sum_i P(ConstraintType_i)=1
\]

donde los tipos disponibles quedan determinados por:

\[
P(Integer)>0 \Rightarrow ARITHMETIC\_LEVEL
\]

\[
P(Real)>0 \Rightarrow ARITHMETIC\_LEVEL
\]

\[
P(String)>0 \Rightarrow TYPE\_LEVEL \land STRING\_CONSTRAINTS
\]


Las distribuciones de operadores deben estar normalizadas. Para operadores booleanos:

\[
P(AND)+P(OR)+P(Implication)+P(Equivalence)=1
\]


Cuando está activo el nivel aritmético, la distribución de operaciones
se valida teniendo en cuenta si las funciones agregadas están
habilitadas. Si están desactivadas:

\[
P(+) + P(-) + P(*) + P(/) = 1
\]

Si las funciones agregadas están activadas, sus probabilidades se
incorporan a la misma distribución:

\[
P(+) + P(-) + P(*) + P(/)
+ P(\mathrm{sum}) + P(\mathrm{avg}) = 1
\]

La distribución de operadores de comparación se valida de forma
separada:

\[
P(\texttt{==}) + P(<) + P(>) + P(\leq) + P(\geq) = 1
\]


\subsubsection*{Validación de atributos}

La configuración de atributos diferencia entre generación aleatoria y definición manual.

Para atributos aleatorios:

\[
0 \leq MIN\_ATTRIBUTES \leq MAX\_ATTRIBUTES
\]

y la distribución de tipos debe cumplir:

\[
P(Boolean)+P(Integer)+P(Real)+P(String)=1
\]


Para atributos definidos manualmente:

\[
MIN\_ATTRIBUTES = NONE
\]

\[
MAX\_ATTRIBUTES = NONE
\]

Además, la utilización de atributos en restricciones depende de los niveles disponibles:

\[
UseInConstraints(Attribute_{Integer|Real})
\Rightarrow Arithmetic\_Level
\]

\[
UseInConstraints(Attribute_{String})
\Rightarrow Type\_Level \land String\_Constraints
\]

En el modo de definición manual, cada elemento de
\texttt{attributes\_list} puede incluir el parámetro
\texttt{attach\_probability}, que determina la probabilidad de asociar el
atributo a una característica. Cuando se proporciona, su valor debe
cumplir:

\[
0 \leq attach\_probability(Attribute) \leq 1
\]

\subsubsection{Consumo por otros módulos}

Una vez validada la configuración, la instancia de 
\texttt{FmgeneratorModel} es utilizada por la operación 
\texttt{GenerateFeatureModel}, encargada de construir el modelo final compatible con FlamaPy.

Esta operación utiliza dicha configuración para generar:

\begin{itemize}
    \item La jerarquía de features.
    \item Los tipos y cardinalidades.
    \item Los atributos asociados.
    \item Las restricciones UVL.
\end{itemize}


La corrección de estas reglas se verifica mediante la suite de pruebas ubicada en:

\begin{lstlisting}
tests/test_validators.py
\end{lstlisting}

donde se incluyen casos que comprueban el rechazo de configuraciones inconsistentes, como distribuciones no normalizadas, dependencias incumplidas o rangos inválidos.

\section{Módulo de generación de artefactos}
\label{sec:generador}

\ph{Describe el componente que produce el artefacto final
(documento JSON, código fuente, fichero CAD, informe PDF\dots) a partir de
la representación intermedia.}

La generación se encuentra encapsulada en la operación 
\texttt{GenerateFeatureModel}, perteneciente a 
\texttt{fm\_generator}. Esta operación recibe como entrada una instancia validada de 
\texttt{FmgeneratorModel}, que contiene todos los parámetros necesarios para controlar la generación, y produce como resultado un objeto 
\texttt{FeatureModel}.

El flujo principal de generación se resume en:

\[
FmgeneratorModel \xrightarrow{\texttt{GenerateFeatureModel}}
FeatureModel
\]

El modelo resultante contiene la estructura jerárquica de features, sus relaciones, atributos asociados y restricciones expresadas mediante el formato interno de FlamaPy, permitiendo posteriormente su transformación a formatos compatibles como UVL.

La ejecución principal se realiza mediante el método:

\begin{lstlisting}[language=Python]
GenerateFeatureModel.execute()
\end{lstlisting}

que valida nuevamente la configuración recibida y coordina las
diferentes fases necesarias para construir el objeto
\texttt{FeatureModel}, cuya transformación a formato UVL se realiza
posteriormente.

\subsection{Cálculos o transformaciones específicas}
\label{subsec:calculos}

\ph{Si la generación implica cálculos no triviales (geometría, fórmulas
analíticas, transformaciones de coordenadas\dots), descríbelos con sus
fórmulas o pseudocódigo.}

La operación \texttt{GenerateFeatureModel} constituye el núcleo del proceso de
generación. Su funcionamiento se divide en varias fases encargadas de
transformar la configuración definida en \texttt{FmgeneratorModel} en un
objeto \texttt{FeatureModel}. Cada fase está encapsulada mediante métodos
independientes que separan la generación de la jerarquía, atributos,
restricciones y configuración específica de UVL.

El algoritmo principal de la operación se muestra en el siguiente
pseudocódigo:

\begin{lstlisting}
FUNCION GenerateFeatureModel.execute:
    INICIO
        VALIDAR configuracion del generador
        INICIALIZAR semilla aleatoria
        feature_model, features =
            GENERAR jerarquia de features
        ASIGNAR atributos a las features generadas
        GENERAR restriciones del modelo
        GENERAR directivas UVL necesarias
        DEVOLVER feature_model
    FIN
\end{lstlisting}

\captionof{lstlisting}{Esquema en pseudocódigo de la operación
\texttt{GenerateFeatureModel.execute}.}

\subsubsection*{Inicialización de la generación}

Antes de comenzar la generación se ejecuta el método
\texttt{\_seed\_generation()}, encargado de configurar la semilla utilizada
por el generador aleatorio.

\begin{lstlisting}
FUNCION _seed_generation(index, attempt):
    INICIO
        CALCULAR seed =
            modelo.seed + index + (attempt * SAT_SEED_STRIDE)
        CONFIGURAR generador aleatorio con seed
    FIN
\end{lstlisting}

\captionof{lstlisting}{Esquema en pseudocódigo del método
\texttt{\_seed\_generation()}.}

\subsubsection*{Generación de la jerarquía de features}

El método
\texttt{\_generate\_hierarchy()} construye la estructura principal del modelo.
Inicialmente crea la feature raíz y determina el número total de features que
serán generadas.

\begin{lstlisting}
FUNCION _generate_hierarchy:
    INICIO
        CREAR feature raiz F0
        DETERMINAR numero_features aleatorio
        CREAR lista de features
        PARA cada feature:
            ASIGNAR tipo booleano o tipado UVL
        PARA cada nivel de profundidad:
            SELECCIONAR features del nivel actual
            SELECCIONAR tipo de relacion padre-hijo
            CREAR relaciones entre features
            ASIGNAR cardinalidades
        DEVOLVER FeatureModel y features generadas
    FIN
\end{lstlisting}

\captionof{lstlisting}{Esquema en pseudocódigo del método
\texttt{\_generate\_hierarchy()}.}

La creación de relaciones depende de las distribuciones configuradas por el
usuario, pudiendo generar relaciones obligatorias, opcionales, alternativas,
disyuntivas o grupos con cardinalidad.

\subsubsection*{Asignación de atributos}

El método \texttt{\_assign\_attributes()} determina el modo de generación de
atributos seleccionado.

\begin{lstlisting}
FUNCION _assign_attributes(features):
    INICIO
        SI modo_aleatorio:
            GENERAR atributos aleatorios
        EN OTRO CASO:
            ASIGNAR atributos definidos manualmente
        ASOCIAR cada atributo con su feature correspondiente
    FIN
\end{lstlisting}

\captionof{lstlisting}{Esquema en pseudocódigo del método
\texttt{\_assign\_attributes()}.}

En el modo aleatorio, el tipo y dominio de cada atributo se obtiene mediante
las distribuciones configuradas. En el modo manual, se utilizan las
definiciones proporcionadas por el usuario durante la configuración del
generador.

\subsubsection*{Generación de restricciones}

La generación de restricciones se realiza mediante el método
\texttt{\_add\_constraints()}, encargado de construir expresiones compatibles
con el metamodelo de FlamaPy.

\begin{lstlisting}
FUNCION _add_constraints(feature_model, features):
    INICIO
        CREAR conjuntos de features y atributos disponibles
        DETERMINAR numero de restricciones
        REPETIR para cada restriccion:
            SELECCIONAR variables compatibles
            DETERMINAR tipo de expresion:
                - booleana
                - aritmetica
                - cadenas
            CONSTRUIR arbol sintactico AST
            ANADIR restriccion al FeatureModel
    FIN
\end{lstlisting}

\captionof{lstlisting}{Esquema en pseudocódigo del método
\texttt{\_add\_constraints()}.}

Cada restricción se representa mediante un árbol sintáctico abstracto
(\texttt{AST}) que permite expresar operaciones lógicas, aritméticas y
funciones específicas de UVL.

\subsubsection*{Generación de directivas UVL}

Finalmente, el método \texttt{\_build\_uvl\_includes()} analiza los niveles
activados y genera las directivas necesarias para representar correctamente
el modelo en UVL.

\begin{lstlisting}
FUNCION _build_uvl_includes:
    INICIO
        CREAR lista de directivas
        SI group_cardinality activo:
            ANADIR Boolean.group-cardinality
        SI feature_cardinality activo:
            ANADIR Arithmetic.feature-cardinality
        SI aggregate_functions activo:
            ANADIR Arithmetic.aggregate-function
        SI string_constraints activo:
            ANADIR Type.string-constraints
        DEVOLVER lista de directivas
    FIN
\end{lstlisting}

\captionof{lstlisting}{Esquema en pseudocódigo del método
\texttt{\_build\_uvl\_includes()}.}

\subsection{Síntesis del artefacto final}
\label{subsec:sintesis}

\ph{Describe el proceso de generación final paso a paso, separando
preparación de datos, renderizado de plantillas y operaciones específicas
de la herramienta destino.}


\subsubsection*{Preparación de la configuración}

Los parámetros definidos mediante el asistente se obtienen del servidor en formato JSON. En el navegador, el archivo \texttt{fmgen\_wrapper.py}, ejecutado dentro de Pyodide, utiliza el método \texttt{from\_flat\_dict} para transformarlos en una instancia de \texttt{FmgeneratorModel}. Antes de iniciar la generación, dicha configuración se valida para comprobar la coherencia entre sus parámetros.


\subsubsection*{Generación del modelo intermedio}

La operación \texttt{GenerateFeatureModel} utiliza la configuración validada
para construir un objeto \texttt{FeatureModel} que contiene la jerarquía de
features, atributos y restricciones generadas.

\[
FmgeneratorModel \rightarrow FeatureModel
\]


\subsubsection*{Transformación a formato UVL}

El modelo generado es transformado mediante las operaciones proporcionadas por
FlamaPy a una representación textual compatible con UVL.

\[
FeatureModel \rightarrow UVL
\]


\subsubsection*{Generación del artefacto final}

Finalmente, la representación UVL obtenida se almacena como fichero de salida,
constituyendo el artefacto final generado por el sistema y listo para ser
utilizado por herramientas compatibles con el lenguaje UVL.

\subsection{Operaciones complementarias}
\label{subsec:operaciones}

\ph{Si tras la generación hay operaciones complementarias
(ensamblado, postprocesado, decoración con metadatos\dots), descríbelas aquí.}

\section{API REST / capa de servicios}
\label{sec:api}

\ph{Describe la API expuesta por el sistema: agrupa los \emph{endpoints} por
bloque temático y, para cada endpoint, indica método, ruta, autenticación
requerida y semántica básica.}

\subsection{Endpoints implementados}
\label{subsec:endpoints}

\ph{Bloques sugeridos:}

\ph{Páginas / vistas:}
\ph{
\begin{itemize}
  \item \ph{\texttt{GET /ruta\_publica}}: \ph{descripción.}
  \item \ph{\texttt{GET /ruta\_protegida}}: \ph{descripción.}
\end{itemize}
}

\ph{Endpoints de lógica principal:}
\ph{
\begin{itemize}
  \item \ph{\texttt{POST /api/\dots}}: \ph{descripción y respuestas (200, 4xx, 5xx).}
  \item \ph{\texttt{GET  /api/\dots}}: \ph{descripción.}
\end{itemize}
}

\ph{Endpoints de autenticación:}
\ph{
\begin{itemize}
  \item \ph{\texttt{POST /api/auth/register}}: \ph{descripción.}
  \item \ph{\texttt{POST /api/auth/login}}: \ph{descripción.}
  \item \ph{\texttt{POST /api/auth/logout}}: \ph{descripción.}
\end{itemize}
}

\ph{Endpoints de gestión de recursos del usuario:}
\ph{
\begin{itemize}
  \item \ph{\texttt{POST /api/\dots}}: \ph{descripción.}
  \item \ph{\texttt{GET  /api/\dots}}: \ph{descripción.}
\end{itemize}
}

\subsubsection*{Endpoints de entrada al generador}

\begin{itemize}
    \item \textbf{GET /generator}: inicializa el asistente de generación,
    reiniciando el estado almacenado en sesión y mostrando la página inicial.

    \item \textbf{GET /generator/random}: redirige al primer paso del flujo de
    generación aleatoria.
\end{itemize}


\subsubsection*{Endpoints del asistente de configuración}

El flujo principal de generación aleatoria se divide en seis etapas. Cada paso
admite peticiones \texttt{GET} para mostrar el formulario correspondiente y
\texttt{POST} para procesar los valores introducidos por el usuario.

\begin{itemize}
    \item \textbf{GET/POST /generator/random/step1}: configuración inicial del
    generador, incluyendo número de modelos y parámetros generales.

    \item \textbf{GET/POST /generator/random/step2}: configuración de los
    niveles de expresividad UVL disponibles.

    \item \textbf{GET/POST /generator/random/step3}: definición de la estructura
    jerárquica del árbol de features.

    \item \textbf{GET/POST /generator/random/step4}: configuración de
    restricciones y operadores utilizados durante la generación.

    \item \textbf{GET/POST /generator/random/step5}: configuración de atributos
    asociados a las features.

    \item \textbf{GET/POST /generator/random/step6}: configuración final,
    generación del modelo y gestión de la salida.
\end{itemize}


\subsubsection*{Endpoints auxiliares}

\begin{itemize}
    \item \textbf{GET /generator/random/params-json}: devuelve la configuración
    actual del asistente en formato JSON.

    \item \textbf{POST /generator/random/summary-refresh/<step>}: actualiza
    dinámicamente el resumen de parámetros mostrado durante la configuración.
\end{itemize}

\subsection{Descarga / exportación de artefactos}
\label{subsec:descarga}

\ph{Si el sistema permite descargar artefactos generados, describe los
escenarios (descarga inmediata tras generación, descarga desde el historial
de configuraciones guardadas\dots) y la estrategia de regeneración bajo demanda
frente a almacenamiento permanente.}

Una vez finalizada la generación, el modelo de características obtenido se
exporta en un formato compatible con UVL para su descarga en ZIP y posterior uso. El sistema no almacena permanentemente los modelos generados, sino que permite
reproducir la generación a partir de la configuración utilizada y la semilla
aleatoria definida.

\section{Integración continua}
\label{sec:ci}

\ph{Describe la automatización de calidad y despliegue del proyecto.
Si usas un proveedor concreto (GitHub Actions, GitLab CI, Jenkins\dots),
nómbralo y comenta los workflows configurados.}

\ph{Workflows típicos:}

\ph{\begin{itemize}
  \item \textbf{\ph{Validación de commits}}: \ph{convención de mensajes, autores\dots}
  \item \textbf{\ph{Linting / análisis estático}}: \ph{herramienta y comandos.}
  \item \textbf{\ph{Pruebas}}: \ph{ejecución de la suite en cada \textit{push} y \textit{PR}.}
  \item \textbf{\ph{Cobertura}}: \ph{umbral mínimo y políticas de fallo.}
  \item \textbf{\ph{Release}}: \ph{condiciones para crear releases.}
  \item \textbf{\ph{Despliegue}}: \ph{construcción y publicación de imagen / artefacto.}
\end{itemize}
}

\ph{Si dispones de comandos locales equivalentes a los workflows remotos
(Make, scripts\dots), enuméralos para que el lector pueda reproducir
las comprobaciones.}

\begin{itemize}
    \item \textbf{Validación de commits}: comprobación automática de los cambios
    introducidos mediante GitHub Actions, verificando que las modificaciones
    cumplen las condiciones establecidas antes de su integración.

    \item \textbf{Linting y análisis estático}: ejecución de herramientas de
    análisis del código para detectar errores de estilo, problemas potenciales
    y mantener una calidad homogénea.

    \item \textbf{Pruebas}: ejecución automática de la suite de tests del
    proyecto en los cambios realizados, asegurando que las nuevas modificaciones
    no afectan a la funcionalidad existente.

    \item \textbf{Despliegue}: automatización de la actualización del entorno de
    ejecución tras superar correctamente las comprobaciones anteriores.
\end{itemize}

% --- Transición al siguiente capítulo ---
\ph{Cierra el capítulo con un párrafo puente de 2--4 líneas: qué queda
establecido aquí y qué pregunta abre el capítulo siguiente. Por ejemplo:
«Con el sistema implementado y la integración continua en marcha, falta
demostrar que se comporta como exigen los requisitos del
Capítulo~\ref{chap:requisitos}. El Capítulo~\ref{chap:pruebas} presenta la
estrategia de pruebas con la que se verifica».}

Con la implementación del generador, la integración de sus componentes y la
automatización de los procesos de validación y despliegue, queda establecido el
funcionamiento completo del sistema desarrollado. El siguiente capítulo presenta
la estrategia de pruebas aplicada, orientada a verificar que la solución cumple
con los requisitos definidos y que los distintos módulos funcionan
correctamente.